\subsection{Howell normal form and canonical remainders}
\label{app:howell}

Howell normal form generalizes reduced row-echelon form
from fields to rings $\Z_m$. It provides a canonical
generating matrix for a submodule and a reduction
procedure for choosing representatives of its cosets.
Throughout, we fix the coordinate order and represent
elements of $\Z_m$ by integers in $\{0,\ldots,m-1\}$.

\begin{definition}[Howell normal form]
    Let $N\leq\Z_m^d$. A matrix $H$ with nonzero rows
    $\vec{h}_1,\ldots,\vec{h}_r$ is a \emph{Howell
    normal form} of $N$ if its rows generate $N$ and:
    \begin{enumerate}
        \item The first nonzero entry of $\vec{h}_i$
        occurs in column $j_i$, with
        $j_1<\cdots<j_r$.

        \item Each pivot $p_i:=H_{i,j_i}$ is a positive
        divisor of $m$.

        \item Entries above a pivot satisfy
        $0\leq H_{\ell,j_i}<p_i$ for $\ell<i$.

        \item For every $j\in\{1,\ldots,d\}$,
        \begin{align}
            \{\vec{a}\in N:
              a_1=\cdots=a_{j-1}=0\}
            =
            \operatorname{span}_{\Z_m}
            \{\vec{h}_i:j_i\geq j\}.
            \label{eq:howell-tail}
        \end{align}
    \end{enumerate}
    Zero rows are omitted.
\end{definition}

The last condition ensures that cancellations in early
coordinates are fully captured by the later rows.
For example, over $\Z_8$, the row $(2,1)$ alone is
in echelon form but not Howell form: multiplying it
by $4$ gives $(0,4)$. The Howell normal form of the
module it generates is
\begin{align}
    H=
    \begin{pmatrix}
        2 & 1\\
        0 & 4
    \end{pmatrix}.
\end{align}

Every submodule has a unique Howell normal form.
Given an $s\times d$ generating matrix, this form
can be computed in polynomial time
\cite{storjohann1998fast}.
Thus it depends only on the embedded submodule,
not on its chosen generators. Since pivot positions
are distinct, it has at most $d$ nonzero rows.

The following lemma gives the reduction procedure
used to choose canonical DCP lifts.

\begin{lemma}[Canonical remainders]
    \label{lem:howell-remainders}
    Let $N\leq\Z_m^d$ have Howell normal form $H$,
    with pivot positions $j_1,\ldots,j_r$ and
    pivots $p_1,\ldots,p_r$.
    Every coset of $N$ contains exactly one vector
    $\vec{a}$ satisfying
    \begin{align}
        0\leq a_{j_i}<p_i
        \qquad (i=1,\ldots,r).
        \label{eq:howell-reduced}
    \end{align}
    This vector can be computed by successively
    reducing the pivot coordinates using the rows
    of $H$.
\end{lemma}

\begin{proof}
    We first construct a reduced representative,
    then prove its uniqueness.

    Start with any representative $\vec{a}$ of
    the coset. At step $i$, write its current
    pivot coordinate as
    \begin{align}
        a_{j_i}=q_ip_i+t_i,
        \qquad 0\leq t_i<p_i,
    \end{align}
    and replace $\vec{a}$ by
    $\vec{a}-q_i\vec{h}_i$, with arithmetic performed
    modulo $m$. This preserves the coset and makes
    the $j_i$-th coordinate equal to $t_i$.
    Since $\vec{h}_i$ vanishes before column $j_i$,
    previously reduced pivot coordinates remain
    unchanged. After all rows have been processed,
    \cref{eq:howell-reduced} holds.

    For uniqueness, suppose $\vec{a}$ and $\vec{b}$
    are distinct reduced representatives of the
    same coset. Then
    $\vec{c}:=\vec{a}-\vec{b}\in N$.
    Let $j$ be the first nonzero coordinate of
    $\vec{c}$. By \cref{eq:howell-tail},
    $\vec{c}$ is a linear combination of rows
    whose pivots occur at or after $j$.
    If $j$ were not a pivot position, all such
    rows would vanish in column $j$, a contradiction.
    Hence $j=j_i$ for some $i$, and $c_{j_i}$
    is a multiple of $p_i$ in $\Z_m$.

    Since $p_i$ divides $m$, the integer
    $a_{j_i}-b_{j_i}$ is therefore divisible by
    $p_i$. But both coordinates lie in
    $\{0,\ldots,p_i-1\}$, so
    \begin{align}
        -(p_i-1)
        \leq a_{j_i}-b_{j_i}
        \leq p_i-1.
    \end{align}
    The only multiple of $p_i$ in this interval
    is zero, contradicting the choice of $j$.
    Thus the reduced representative is unique.
\end{proof}

Write $\operatorname{rem}_H(\vec{a})$ for this
unique representative of $\vec{a}+N$.
The lemma immediately gives
\begin{align}
    \vec{a}\in N
    &\iff \operatorname{rem}_H(\vec{a})=\vec{0},
    \\
    \vec{a}+N=\vec{b}+N
    &\iff
    \operatorname{rem}_H(\vec{a})
    =\operatorname{rem}_H(\vec{b}).
\end{align}
Once $H$ is known, reduction takes $O(rd)$
coordinate operations and integer divisions on
representatives smaller than $m$.
In particular, for the fixed modulus $m=8$, it
takes $O(d^2)$ time. The remainder map need not
be linear; it selects representatives of quotient
classes without asserting that the quotient splits.

To apply this construction to a DCP subgroup $G$,
take
\begin{align}
    N=\iota(\Delta(G))\leq\Z_8^d,
    \qquad d=1+n+\binom{n}{2},
\end{align}
and let $H$ be its Howell normal form.
For a lift $h_j=D_jX(\vec{x}^{(j)})\in G$, define
$\widetilde{D}_j$ by
\begin{align}
    \iota(\widetilde{D}_j)
    =
    \operatorname{rem}_H\big(\iota(D_j)\big).
\end{align}
This is a valid diagonal DCP operator: reduction
changes $\iota(D_j)$ by an element of
$\iota(\Delta(G))$, so the result remains in the
image of $\iota$. Moreover,
$\widetilde{D}_jD_j^{-1}\in\Delta(G)$, and hence
\begin{align}
    \widetilde{h}_j
    :=
    \widetilde{D}_jX(\vec{x}^{(j)})
    =
    (\widetilde{D}_jD_j^{-1})h_j
    \in G.
\end{align}
All lifts of $\vec{x}^{(j)}$ have diagonal coordinate
vectors in the same coset of $N$. Therefore
$\widetilde{h}_j$ depends only on $G$, $\vec{x}^{(j)}$,
and the fixed coordinate order, not on the
initial lift.
